\documentclass[tikz,border=8pt]{standalone}
\usepackage{tikz}
\usepackage{amsmath}
\usepackage{amssymb}
\usepackage{pgfplots}
\pgfplotsset{compat=1.18}
\usepackage{ctex}
\usetikzlibrary{calc,positioning,arrows.meta,matrix}

% LC 53 Maximum Subarray (Kadane)
% 数组 [-2,1,-3,4,-1,2,1,-5,4]
% 追踪 cur_max / global_max 变化，最终 global_max=6，子数组[4,-1,2,1]

\begin{document}
\begin{tikzpicture}[
  every node/.style={font=\small},
  cell/.style={draw=gray!60, rectangle,
    minimum width=0.8cm, minimum height=0.65cm, thick, fill=white},
  hilightcell/.style={draw=red!60, rectangle,
    minimum width=0.8cm, minimum height=0.65cm, very thick, fill=red!10},
]

%% ────────────────────────────────────────────────────────
%% 结构（从上到下）：
%%  pgfplots chart（顶部，origin at scope shift=(0,0)）
%%  标题/数组/表格/说明框 均在 chart 下方
%%
%% pgfplots height=6cm；chart TikZ box 高度 ≈ 6cm
%% chart bottom ≈ y = -6 after shift 0
%% 标题 at y = -8, 数组 at y=-9.2
%% ────────────────────────────────────────────────────────

%% ── pgfplots 折线图（顶部，从 y=0 开始向下）────────────
\begin{axis}[
  width=13cm,
  height=6cm,
  xmin=-0.5, xmax=8.5,
  ymin=-3.5, ymax=7.0,
  xlabel={索引 $i$},
  ylabel={值},
  xlabel style={font=\small},
  ylabel style={font=\small},
  xtick={0,1,2,3,4,5,6,7,8},
  xticklabels={0,1,2,3,4,5,6,7,8},
  ytick={-3,-2,-1,0,1,2,3,4,5,6,7},
  tick label style={font=\scriptsize},
  grid=both,
  grid style={gray!20,thin},
  major grid style={gray!30},
  axis line style={gray!60},
  legend pos=north west,
  legend style={font=\scriptsize,fill=white,draw=gray!40},
  title={Kadane 最大子数组（LC 53 Maximum Subarray）},
  title style={font=\large\bfseries,yshift=4pt},
]

% cur_max 折线
\addplot[orange!70!red, thick, mark=*, mark size=2.5pt,
         mark options={fill=orange!50}]
  coordinates {
    (0,-2)(1,1)(2,-2)(3,4)(4,3)(5,5)(6,6)(7,1)(8,5)
  };
\addlegendentry{cur\_max};

% global_max 折线
\addplot[blue!60, very thick, mark=square*, mark size=2.5pt,
         mark options={fill=blue!30}]
  coordinates {
    (0,-2)(1,1)(2,1)(3,4)(4,4)(5,5)(6,6)(7,6)(8,6)
  };
\addlegendentry{global\_max};

% global_max=6 水平线
\addplot[red!50, dashed, thick] coordinates {(-0.5,6)(8.5,6)};

% 峰值标注
\node[font=\scriptsize,red!70,anchor=south west] at (axis cs:6.1,6.05)
  {global\_max=6};

% 子数组范围高亮区域（idx 3..6）
\fill[red!8,opacity=0.5] (axis cs:2.5,-3.5) rectangle (axis cs:6.5,7.0);
\draw[red!40,dashed,thin] (axis cs:2.5,-3.5) -- (axis cs:2.5,7.0);
\draw[red!40,dashed,thin] (axis cs:6.5,-3.5) -- (axis cs:6.5,7.0);
\node[font=\tiny,red!60,anchor=north] at (axis cs:4.5,-3.3)
  {子数组[3..6]};

\end{axis}

%% ── 原始数组（chart 下方）───────────────────────────────
\node[font=\small,gray] at (7.5, -7.8)
  {原始数组：\texttt{[-2, 1, -3, 4, -1, 2, 1, -5, 4]}};

\node[font=\scriptsize,gray,anchor=east] at (0.5, -8.6) {数组：};
\foreach \i/\v in {0/{-2}, 1/1, 2/{-3}, 3/4, 4/{-1}, 5/2, 6/1, 7/{-5}, 8/4} {
  \pgfmathsetmacro\xi{1.0+\i*1.35}
  \ifnum\i>2
    \ifnum\i<7
      \node[hilightcell,font=\ttfamily\small] at (\xi, -8.6) {\v};
    \else
      \node[cell,font=\ttfamily\small] at (\xi, -8.6) {\v};
    \fi
  \else
    \node[cell,font=\ttfamily\small] at (\xi, -8.6) {\v};
  \fi
  \node[font=\tiny,gray] at (\xi, -9.15) {\i};
}
\node[font=\scriptsize,red!70] at (5.35, -9.7)
  {红色高亮 = 最大子数组 [4,\,-1,\,2,\,1]，和=6};

%% ── 逐步数值表 ────────────────────────────────────────────
\node[font=\scriptsize,gray,anchor=west] at (0.5, -10.5) {步骤明细：};

\node[font=\scriptsize\bfseries,gray]      at (1.6,  -10.9) {idx};
\node[font=\scriptsize\bfseries,gray]      at (3.0,  -10.9) {nums[i]};
\node[font=\scriptsize\bfseries,orange!70] at (5.0,  -10.9) {cur\_max};
\node[font=\scriptsize\bfseries,blue!60]   at (7.2,  -10.9) {global\_max};
\node[font=\scriptsize\bfseries,gray]      at (10.2, -10.9) {操作};

\foreach \i/\v/\cm/\gm/\op in {
  0/{-2}/{-2}/{-2}/{初始化},
  1/1/1/1/{重置(1>-2+1)},
  2/{-3}/{-2}/1/{延续(-2>-3)},
  3/4/4/4/{重置(4>-2+4)},
  4/{-1}/3/4/{延续(3>-1)},
  5/2/5/5/{延续(5>3+2)},
  6/1/6/6/{延续(6>5+1) $\to$ 峰值},
  7/{-5}/1/6/{延续(1>6-5)},
  8/4/5/6/{延续(5>1+4)}} {
  \pgfmathsetmacro\ty{-11.4-\i*0.55}
  \node[font=\scriptsize] at (1.6, \ty) {\i};
  \node[font=\scriptsize,font=\ttfamily] at (3.0, \ty) {\v};
  \node[font=\scriptsize,orange!70] at (5.0, \ty) {\cm};
  \node[font=\scriptsize,blue!60] at (7.2, \ty) {\gm};
  \node[font=\scriptsize,gray,anchor=west] at (8.2, \ty) {\op};
}

%% ── 分隔线 ──────────────────────────────────────────────
\draw[gray!25,dashed] (0.0, -17.1) -- (15.0, -17.1);

%% ── 算法说明 ─────────────────────────────────────────────
\node[draw=gray!50,fill=white,rounded corners=3pt,font=\small,
      inner sep=8pt,align=left]
  at (7.5, -18.1)
  {\textbf{Kadane 算法（$O(n)$ 时间，$O(1)$ 空间）：}\\[4pt]
   \texttt{cur\_max = global\_max = nums[0]}\\[2pt]
   遍历 \texttt{i = 1..n-1}：\\
   \quad \texttt{cur\_max = max(nums[i], cur\_max + nums[i])}（当前子数组要或不要前缀）\\
   \quad \texttt{global\_max = max(global\_max, cur\_max)}};

\end{tikzpicture}
\end{document}
